\scanpage{4}
\begin{center}\textbf{Introduction}\end{center}

Soit $K$ un corps local de caractéristique 0 et de caractéristique résiduelle $p$, et soit $C$ la complétion d'une clôture algébrique de $K$. Soit $T$ le module de Tate ([9], n\textsuperscript{o} 2.4) associé à un groupe $p$-divisible $F$, défini sur l'anneau des entiers de $K$. \textsc{Tate} a montré ([9], \S\ 4, cor. 2 au th. 3) que $T\otimes C$ possède une décomposition analogue à celle de Hodge pour la cohomologie complexe. De nombreuses propriétés du module galoisien $T$ sont implicitement contenues dans cette décomposition. Dans son séminaire au Collège de France, résumé dans [8], \textsc{Tate} en a indiqué un certain nombre (notamment lorsque l'action du groupe de Galois est abélienne). Dans ce qui suit, j'explicite une autre conséquence de cette décomposition de $T\otimes C$: si $G$ est le groupe de Galois qui opère sur $T$, l'enveloppe algébrique de $G$ contient (sous une hypothèse de semi-simplicité convenable) un groupe «de Mumford-Tate» $p$-adique (\S\ 3, ths. 1 et 2). Lorsqu'on fait certaines hypothèses supplémentaires sur $F$, on en déduit que $G$ est \textit{ouvert} dans $\operatorname{Aut}(T)$ (cf. \S\ 5, th. 4). Ces hypothèses sont notamment vérifiées lorsque $F$ est un groupe formel à 1 paramètre, n'admettant pas de multiplication complexe formelle (\S\ 5, th. 5).

\begin{center}\textbf{\S\ 1. Enveloppes algébriques des groupes\par linéaires $p$-adiques}\end{center}

Soit $V$ un espace vectoriel de dimension finie sur le corps $p$-adique $\mathbf Q_p$, et soit $\operatorname{Aut}(V)$ le groupe de ses automorphismes, muni de sa structure

\footnotetext[1]{Le texte ci-dessous a été rédigé en décembre 1966. Il diffère de l'exposé oral par un plus grand usage de la théorie des groupes algébriques.

L'idée de remplacer les algèbres de Lie $p$-adiques par les groupes algébriques correspondants m'avait d'ailleurs été suggérée par Grothendieck il y a plusieurs années, en liaison avec sa théorie des «motifs».}
\margin{3}
